\documentclass[10pt,a4paper]{amsart}
\usepackage{amsmath,amssymb,amsthm,graphicx}
\usepackage{verbatim}
\newenvironment{reference}{\comment}{\endcomment}
%\newenvironment{reference}{}{}
\newenvironment{slogan}{\comment}{\endcomment}
\newenvironment{history}{\comment}{\endcomment}
% For commutative diagrams we use Xy-pic
\usepackage[all]{xy}

% We use 2cell for 2-commutative diagrams.
\xyoption{2cell}
% Theorem environments.
%
\theoremstyle{plain}
\newtheorem{theorem}[subsection]{Theorem}
\newtheorem{proposition}[subsection]{Proposition}
\newtheorem{lemma}[subsection]{Lemma}

\theoremstyle{definition}
\newtheorem{definition}[subsection]{Definition}
\newtheorem{example}[subsection]{Example}
\newtheorem{exercise}[subsection]{Exercise}
\newtheorem{situation}[subsection]{Situation}

\theoremstyle{remark}
\newtheorem{remark}[subsection]{Remark}
\newtheorem{remarks}[subsection]{Remarks}

\numberwithin{equation}{subsection}

% Macros
%
\def\lim{\mathop{\mathrm{lim}}\nolimits}
\def\colim{\mathop{\mathrm{colim}}\nolimits}
\def\Spec{\mathop{\mathrm{Spec}}}
\def\Hom{\mathop{\mathrm{Hom}}\nolimits}
\def\Ext{\mathop{\mathrm{Ext}}\nolimits}
\def\SheafHom{\mathop{\mathcal{H}\!\mathit{om}}\nolimits}
\def\SheafExt{\mathop{\mathcal{E}\!\mathit{xt}}\nolimits}
\def\Sch{\mathit{Sch}}
\def\Mor{\mathop{\mathrm{Mor}}\nolimits}
\def\Ob{\mathop{\mathrm{Ob}}\nolimits}
\def\Sh{\mathop{\mathit{Sh}}\nolimits}
\def\NL{\mathop{N\!L}\nolimits}
\def\CH{\mathop{\mathrm{CH}}\nolimits}
\def\proetale{{pro\text{-}\acute{e}tale}}
\def\etale{{\acute{e}tale}}
\def\QCoh{\mathit{QCoh}}
\def\Ker{\mathop{\mathrm{Ker}}}
\def\Im{\mathop{\mathrm{Im}}}
\def\Coker{\mathop{\mathrm{Coker}}}
\def\Coim{\mathop{\mathrm{Coim}}}

% Boxtimes
%
\DeclareMathSymbol{\boxtimes}{\mathbin}{AMSa}{"02}

%
% Macros for moduli stacks/spaces
%
\def\QCohstack{\mathcal{QC}\!\mathit{oh}}
\def\Cohstack{\mathcal{C}\!\mathit{oh}}
\def\Spacesstack{\mathcal{S}\!\mathit{paces}}
\def\Quotfunctor{\mathrm{Quot}}
\def\Hilbfunctor{\mathrm{Hilb}}
\def\Curvesstack{\mathcal{C}\!\mathit{urves}}
\def\Polarizedstack{\mathcal{P}\!\mathit{olarized}}
\def\Complexesstack{\mathcal{C}\!\mathit{omplexes}}
% \Pic is the operator that assigns to X its picard group, usage \Pic(X)
% \Picardstack_{X/B} denotes the Picard stack of X over B
% \Picardfunctor_{X/B} denotes the Picard functor of X over B
\def\Pic{\mathop{\mathrm{Pic}}\nolimits}
\def\Picardstack{\mathcal{P}\!\mathit{ic}}
\def\Picardfunctor{\mathrm{Pic}}
\def\Deformationcategory{\mathcal{D}\!\mathit{ef}}
\usepackage{cleveref}
\crefname{originalreference}{original reference}{original references}
\Crefname{originalreference}{Original reference}{Original references}
\makeatletter
\newlabel{coherent-lemma-property-higher-rank-cohomological}{{coherent-lemma-property-higher-rank-cohomological}{source}{}{sourcepointer-1}{}}
\newlabel{coherent-lemma-property-higher-rank-cohomological@cref}{{[originalreference][0][]coherent-lemma-property-higher-rank-cohomological}{[1][0][]source}{}{}{}}
\makeatother
\title{Vanishing of the top local cohomology of one finite module of full support forces the cohomological dimension of the ideal to drop}
\author{The Stacks Project Authors}
\date{stacks-project:local-cohomology.tex:local-snapshot-20261012; GNU FDL 1.2 or later}
\begin{document}
\maketitle
\noindent\textit{Editorial note.} All delivered mathematical text is the author's own, copied from the pinned Stacks Project chapter source sources/stacks-local-cohomology/pinned-original-source.tex; no statement or proof step was rewritten, reordered or removed.\par\medskip
\noindent\textit{Editorial note.} Preamble provenance: the Stacks Project keeps one single shared preamble for every chapter; this card takes it directly from the staged sources/stacks-local-cohomology/preamble.tex as three verbatim spans (lines 8-12, 14-18 and 154-217), so no derived file was written and nothing was added or redefined. The wrapper class and the base packages are supplied by the assembler.\par\medskip
\noindent\textit{Editorial note.} Bibliography provenance: neither the counted statement nor its proof contains a citation (this lane rescans the whole span before writing the card), so the card carries no bibliography and no entry was reconstructed.\par\medskip
\noindent\textit{Editorial note.} Selection provenance: this card was opened under the file-based control contract (CONTROL.md), with the mechanical pre-card gate run on its statement and proof interval before the card was written.\par\medskip
\noindent\textit{Allowed standard prerequisites (editorial).}\begin{itemize}
\item The vanishing of local cohomology above the dimension of the support and the bound cd(A, I) \textbackslash{}leq dim(A). Author lemma lemma-cd-dimension, whose statement is reproduced verbatim as a context span of this card (source lines 758-764). Gives the a priori bound cd(A, I) \textbackslash{}leq d and the vanishing that makes the property P descend along surjections.
\item The statement that a property of finite modules which holds for a module of full support and is inherited by quotients of direct sums holds for every finite module. Invoked at the end of the counted proof (source lines 3626-3628) and attributed to a lemma of another chapter (a cross-chapter pointer). Spreads the vanishing from the given module to all finite modules, which is what cd(A, I) < d means.
\end{itemize}\medskip
\ifdefined\hypertarget\hypertarget{sourcepointer-1}{}\fi
\noindent\textit{Source reference.} Original reference coherent-lemma-property-higher-rank-cohomological: sources/stacks-local-cohomology/pinned-original-source.tex, source lines 3627--3627. Navigation pointer only: the counted text refers to the Stacks item coherent-lemma-property-higher-rank-cohomological, which belongs to another chapter of the project and is not part of this card. The pointer carries the target's own label as its number because this chapter's staged snapshot does not contain the other chapter's tag table; its content is not reproduced here.\par\medskip
\begin{lemma}
\label{lemma-cd-top-vanishing-some-module}
Let $A$ be a Noetherian ring of dimension $d$. Let $I \subset A$
be an ideal. If $H^d_{V(I)}(M) = 0$ for some finite $A$-module
whose support contains all the irreducible components of
dimension $d$, then $\text{cd}(A, I) < d$.
\end{lemma}

\begin{lemma}
\label{lemma-cd-dimension}
Let $I \subset A$ be a finitely generated ideal of a ring $A$.
If $M$ is a finite $A$-module, then
$H^i_{V(I)}(M) = 0$ for $i > \dim(\text{Supp}(M))$.
In particular, we have $\text{cd}(A, I) \leq \dim(A)$.
\end{lemma}

\begin{proof}
By Lemma \ref{lemma-cd-dimension} we know $\text{cd}(A, I) \leq d$.
Thus for any finite $A$-module $N$ we have $H^i_{V(I)}(N) = 0$
for $i > d$. Let us say property $\mathcal{P}$ holds for the
finite $A$-module $N$ if $H^d_{V(I)}(N) = 0$.
One of our assumptions is that $\mathcal{P}(M)$ holds.
Observe that $\mathcal{P}(N_1 \oplus N_2)
\Leftrightarrow (\mathcal{P}(N_1) \wedge \mathcal{P}(N_2))$.
Observe that if $N \to N'$ is surjective, then
$\mathcal{P}(N) \Rightarrow \mathcal{P}(N')$ as we
have the vanishing of $H^{d + 1}_{V(I)}$ (see above).
Let $\mathfrak p_1, \ldots, \mathfrak p_n$ be the
minimal primes of $A$ with $\dim(A/\mathfrak p_i) = d$.
Observe that $\mathcal{P}(N)$ holds if the support
of $N$ is disjoint from $\{\mathfrak p_1, \ldots, \mathfrak p_n\}$
for dimension reasons, see Lemma \ref{lemma-cd-dimension}.
For each $i$ set $M_i = M/\mathfrak p_i M$.
This is a finite $A$-module annihilated by $\mathfrak p_i$
whose support is equal to
$V(\mathfrak p_i)$ (here we use the assumption on the support of $M$).
Finally, if $J \subset A$ is an ideal, then we have $\mathcal{P}(JM_i)$
as $JM_i$ is a quotient of a direct sum of copies of $M$.
Thus it follows from Cohomology of Schemes, Lemma
\ref{coherent-lemma-property-higher-rank-cohomological}
that $\mathcal{P}$ holds for every finite $A$-module.
\end{proof}

\end{document}
